\documentclass[11pt,a4paper]{article}
\usepackage{fontspec,amsmath,amssymb,amsthm,microtype}
\setmainfont{lmroman10-regular.otf}[BoldFont=lmroman10-bold.otf,ItalicFont=lmroman10-italic.otf,BoldItalicFont=lmroman10-bolditalic.otf]
\usepackage{xcolor}
\usepackage[most]{tcolorbox}
\usepackage{mathrsfs}
\providecommand{\lt}{<}
\providecommand{\gt}{>}
\setmonofont{DejaVuSansMono.ttf}[BoldFont=DejaVuSansMono-Bold.ttf,ItalicFont=DejaVuSansMono-Oblique.ttf,BoldItalicFont=DejaVuSansMono-BoldOblique.ttf,Scale=0.86]
\newfontfamily\leanfallback{STIXTwoMath-Regular.otf}[Scale=0.86]
\definecolor{leanrule}{gray}{0.62}
\definecolor{leanback}{gray}{0.962}
\definecolor{leantitle}{gray}{0.22}
\definecolor{leanmuted}{gray}{0.38}
\newtcolorbox{leanbox}[1]{enhanced,breakable,colback=leanback,colframe=leanrule,
  boxrule=0.4pt,arc=0.6mm,left=2.4mm,right=2.4mm,top=1.8mm,bottom=1.8mm,
  before skip=0.8em,after skip=1em,fonttitle=\small\bfseries,coltitle=leantitle,
  colbacktitle=leanback,titlerule=0.3pt,title={#1}}
\newlength{\leancharwidth}
\newenvironment{leancode}{\par\smallskip\begingroup\ttfamily\footnotesize\raggedright
  \settowidth{\leancharwidth}{M}\setlength{\parskip}{0pt}\setlength{\parindent}{0pt}
  \frenchspacing}{\par\endgroup\smallskip}
\newcommand{\leanline}[2]{\hangindent=\dimexpr#1\leancharwidth+4\leancharwidth\relax
  \hangafter=1\noindent\hspace*{#1\leancharwidth}#2\par}
\newcommand{\leansource}[1]{\par{\footnotesize\color{leanmuted}#1}\par}
\newcommand{\leanlabel}[1]{\par\smallskip\noindent{\small\bfseries #1}\par}
\usepackage[margin=25mm]{geometry}
\usepackage[unicode,hidelinks]{hyperref}
\hypersetup{pdfauthor={LangevinLean Project}}
\setlength{\emergencystretch}{3em}
\setlength{\parskip}{0.4em}
\setlength{\parindent}{0pt}
\allowdisplaybreaks[2]
\raggedbottom
\title{\texorpdfstring{A global actual Gibbs diffusion and its continuous operators}{A global actual Gibbs diffusion and its continuous operators}}
\author{LangevinLean Project}
\date{}
\begin{document}\maketitle
Draft mathematical exposition; these notes may contain mistakes.


Finite-horizon Langevin solutions driven by the same actual Brownian source agree on overlaps by uniqueness. Their endpoints define a global canonical path. It has the Gibbs law at every nonnegative time and yields strongly continuous actual observable contractions.


\section{\texorpdfstring{Actual Brownian input and finite-horizon solutions}{Actual Brownian input and finite-horizon solutions}}


\textbf{Construction.} Use the coordinate Brownian source on the independent Gibbs extension and the actual Gibbs initial variable. For each finite horizon, Picard existence gives a continuous \(L^2\) solution of

\[Y_t-Y_0=-\int_0^t\nabla V(Y_s)\,ds+\sqrt2 B_t.\]

The initial law is the actual normalized Gibbs measure; Brownian coordinates retain their actual pre-Brownian laws on the product space.

\begin{leanbox}{Lean: \texttt{brownian}, \texttt{brownian\_\allowbreak{}preBrownian}, \texttt{initial} and 3 more}
\leansource{\texttt{brownian} \quad WassersteinMixing/\allowbreak{}SemigroupGlobal.lean, lines 20–21}
\begin{leancode}
\leanline{0}{def\ brownian\ (P\ :\ Parameters)\ (i\ :\ Fin\ P.dimension)\ (t\ :\ ℝ≥0)\ (ω\ :\ ActualSpace\ P)\ :\ ℝ\ :=}
\leanline{2}{ULA.CanonicalBrownian.coordinate\ i\ t\ ω.2}
\end{leancode}
\leansource{\texttt{brownian\_\allowbreak{}preBrownian} \quad WassersteinMixing/\allowbreak{}SemigroupGlobal.lean, lines 23–29}
\begin{leancode}
\leanline{0}{theorem\ brownian\_preBrownian\ (P\ :\ Parameters)\ (V\ :\ E\ P.dimension\ →\ ℝ)}
\leanline{4}{(hV\ :\ PotentialAssumptions\ P\ V)\ :}
\leanline{4}{∀\ i,\ IsPreBrownianReal\ (brownian\ P\ i)\ (actualProbability\ P\ V)}
\end{leancode}
\leansource{\texttt{initial} \quad WassersteinMixing/\allowbreak{}SemigroupGlobal.lean, lines 32–36}
\begin{leancode}
\leanline{0}{def\ initial\ (P\ :\ Parameters)\ (V\ :\ E\ P.dimension\ →\ ℝ)\ (hV\ :\ PotentialAssumptions\ P\ V)\ :}
\leanline{4}{Lp\ (E\ P.dimension)\ 2\ (actualProbability\ P\ V)\ :=\ by}
\leanline{2}{let\ :=\ actualTarget\_isProbability\ P\ V\ hV}
\leanline{2}{exact\ ULA.initialExtension\_lpReference\ (diagonalGibbsCoupling\ P\ V)}
\leanline{4}{(ULA.CanonicalBrownian.vectorProbability\ (Fin\ P.dimension))\ (actualTarget\_memLp\ P\ V\ hV)}
\end{leancode}
\leansource{\texttt{initial\_\allowbreak{}law} \quad WassersteinMixing/\allowbreak{}SemigroupGlobal.lean, lines 38–43}
\begin{leancode}
\leanline{0}{theorem\ initial\_law\ (P\ :\ Parameters)\ (V\ :\ E\ P.dimension\ →\ ℝ)}
\leanline{4}{(hV\ :\ PotentialAssumptions\ P\ V)\ :}
\leanline{4}{(actualProbability\ P\ V).map\ (initial\ P\ V\ hV)\ =\ ULA.gibbsMeasure\ V}
\end{leancode}
\leansource{\texttt{solutionOn} \quad WassersteinMixing/\allowbreak{}SemigroupGlobal.lean, lines 46–51}
\begin{leancode}
\leanline{0}{def\ solutionOn\ (P\ :\ Parameters)\ (V\ :\ E\ P.dimension\ →\ ℝ)\ (hV\ :\ PotentialAssumptions\ P\ V)}
\leanline{4}{(T\ :\ ℝ)\ (hT\ :\ 0\ ≤\ T)\ :\ ℝ\ →\ Lp\ (E\ P.dimension)\ 2\ (actualProbability\ P\ V)\ :=\ by}
\leanline{2}{let\ :=\ actualProbability\_isProbability\ P\ V\ hV}
\leanline{2}{exact\ Classical.choose\ (ULA.exists\_lp\_langevin\_integral\_of\_brownian\ (actualProbability\ P\ V)}
\leanline{4}{V\ ⟨P.beta,\ P.beta\_pos.le⟩\ hV.gradientLipschitz\ (initial\ P\ V\ hV)}
\leanline{4}{(brownian\ P)\ (brownian\_preBrownian\ P\ V\ hV)\ 0\ T\ le\_rfl\ hT)}
\end{leancode}
\leansource{\texttt{solutionOn\_\allowbreak{}spec} \quad WassersteinMixing/\allowbreak{}SemigroupGlobal.lean, lines 53–66}
\begin{leancode}
\leanline{0}{theorem\ solutionOn\_spec\ (P\ :\ Parameters)\ (V\ :\ E\ P.dimension\ →\ ℝ)}
\leanline{4}{(hV\ :\ PotentialAssumptions\ P\ V)\ (T\ :\ ℝ)\ (hT\ :\ 0\ ≤\ T)\ :}
\leanline{4}{letI\ :=\ actualProbability\_isProbability\ P\ V\ hV}
\leanline{4}{Continuous\ (solutionOn\ P\ V\ hV\ T\ hT)\ ∧\ solutionOn\ P\ V\ hV\ T\ hT\ 0\ =\ initial\ P\ V\ hV\ ∧}
\leanline{4}{(∀\ t\ ∈\ Icc\ 0\ T,\ solutionOn\ P\ V\ hV\ T\ hT\ t\ -\ initial\ P\ V\ hV\ =}
\leanline{6}{-(∫\ s\ in\ 0..t,\ ULA.lpDrift\ (actualProbability\ P\ V)\ (gradient\ V)}
\leanline{8}{hV.gradientLipschitz\ (solutionOn\ P\ V\ hV\ T\ hT\ s))\ +}
\leanline{6}{ULA.brownianLangevinNoise\ (actualProbability\ P\ V)\ (brownian\ P)}
\leanline{8}{(brownian\_preBrownian\ P\ V\ hV)\ 0\ t)}
\end{leancode}
\leanlabel{How the Lean proof follows the argument}
\texttt{brownian} and \texttt{initial} select the actual source and initialization. \texttt{solution\allowbreak{}On} chooses a solution only from the proved existence theorem. \texttt{solution\allowbreak{}On\_\allowbreak{}spec} retains continuity, the exact initial value and the actual integral equation.
\end{leanbox}


Earlier results used here: \cite[Section~3.13.2]{LangevinLean2026ULASC}; \cite[Section~3.22.20]{LangevinLean2026ULASC}; \cite[Section~3.22.22]{LangevinLean2026ULASC}; \cite[Section~3.25.1]{LangevinLean2026ULASC}; \cite[Section~3.25.6]{LangevinLean2026ULASC}; \cite[Section~3.26.3]{LangevinLean2026ULASC}; \cite[Section~3.32.12]{LangevinLean2026ULASC}; \cite[Section~3.8.17]{LangevinLean2026ULASC}; \cite[Section~3.8.19]{LangevinLean2026ULASC}; \cite[Section~3.8.20]{LangevinLean2026ULASC}.

\section{\texorpdfstring{Consistency, continuity and identification with the reference path}{Consistency, continuity and identification with the reference path}}


\textbf{Claim.} Define the global endpoint at \(t\ge0\) using horizon \(t\), and hold it constant before zero. It agrees with every longer-horizon solution, is continuous, satisfies the same integral equation and coincides with the actual comparison segment and grid reference.



\textbf{Proof.} Uniqueness of solutions with the same initial state and Brownian input identifies overlapping finite horizons. Near any time choose one fixed larger horizon; the global path equals that continuous solution locally, including the clipped-zero presentation. Substituting consistency into the finite-horizon integral equation gives the global one. The actual segmented reference solves that same equation, so uniqueness gives its segment and grid identification.

\begin{leanbox}{Lean: \texttt{state}, \texttt{state\_\allowbreak{}eq\_\allowbreak{}solutionOn}, \texttt{state\_\allowbreak{}zero} and 6 more}
\leansource{\texttt{state} \quad WassersteinMixing/\allowbreak{}SemigroupGlobal.lean, lines 70–72}
\begin{leancode}
\leanline{0}{def\ state\ (P\ :\ Parameters)\ (V\ :\ E\ P.dimension\ →\ ℝ)\ (hV\ :\ PotentialAssumptions\ P\ V)}
\leanline{4}{(t\ :\ ℝ)\ :\ Lp\ (E\ P.dimension)\ 2\ (actualProbability\ P\ V)\ :=}
\leanline{2}{solutionOn\ P\ V\ hV\ (max\ t\ 0)\ (le\_max\_right\ \_\ \_)\ (max\ t\ 0)}
\end{leancode}
\leansource{\texttt{state\_\allowbreak{}eq\_\allowbreak{}solutionOn} \quad WassersteinMixing/\allowbreak{}SemigroupGlobal.lean, lines 75–88}
\begin{leancode}
\leanline{0}{theorem\ state\_eq\_solutionOn\ (P\ :\ Parameters)\ (V\ :\ E\ P.dimension\ →\ ℝ)}
\leanline{4}{(hV\ :\ PotentialAssumptions\ P\ V)\ (T\ :\ ℝ)\ (hT\ :\ 0\ ≤\ T)}
\leanline{4}{(t\ :\ ℝ)\ (ht\ :\ 0\ ≤\ t)\ (htT\ :\ t\ ≤\ T)\ :}
\leanline{4}{state\ P\ V\ hV\ t\ =\ solutionOn\ P\ V\ hV\ T\ hT\ t}
\end{leancode}
\leansource{\texttt{state\_\allowbreak{}zero} \quad WassersteinMixing/\allowbreak{}SemigroupGlobal.lean, lines 90–92}
\begin{leancode}
\leanline{0}{@[simp]\ theorem\ state\_zero\ (P\ :\ Parameters)\ (V\ :\ E\ P.dimension\ →\ ℝ)}
\leanline{4}{(hV\ :\ PotentialAssumptions\ P\ V)\ :\ state\ P\ V\ hV\ 0\ =\ initial\ P\ V\ hV}
\end{leancode}
\leansource{\texttt{state\_\allowbreak{}eq\_\allowbreak{}solutionOn\_\allowbreak{}max} \quad WassersteinMixing/\allowbreak{}SemigroupGlobal.lean, lines 94–99}
\begin{leancode}
\leanline{0}{theorem\ state\_eq\_solutionOn\_max\ (P\ :\ Parameters)\ (V\ :\ E\ P.dimension\ →\ ℝ)}
\leanline{4}{(hV\ :\ PotentialAssumptions\ P\ V)\ (T\ :\ ℝ)\ (hT\ :\ 0\ ≤\ T)}
\leanline{4}{(t\ :\ ℝ)\ (htT\ :\ max\ t\ 0\ ≤\ T)\ :}
\leanline{4}{state\ P\ V\ hV\ t\ =\ solutionOn\ P\ V\ hV\ T\ hT\ (max\ t\ 0)}
\end{leancode}
\leansource{\texttt{continuous\_\allowbreak{}state} \quad WassersteinMixing/\allowbreak{}SemigroupGlobal.lean, lines 103–115}
\begin{leancode}
\leanline{0}{theorem\ continuous\_state\ (P\ :\ Parameters)\ (V\ :\ E\ P.dimension\ →\ ℝ)}
\leanline{4}{(hV\ :\ PotentialAssumptions\ P\ V)\ :\ Continuous\ (state\ P\ V\ hV)}
\end{leancode}
\leansource{\texttt{continuousOn\_\allowbreak{}state} \quad WassersteinMixing/\allowbreak{}SemigroupGlobal.lean, lines 118–130}
\begin{leancode}
\leanline{0}{theorem\ continuousOn\_state\ (P\ :\ Parameters)\ (V\ :\ E\ P.dimension\ →\ ℝ)}
\leanline{4}{(hV\ :\ PotentialAssumptions\ P\ V)\ :\ ContinuousOn\ (state\ P\ V\ hV)\ (Ici\ 0)}
\end{leancode}
\leansource{\texttt{state\_\allowbreak{}integral} \quad WassersteinMixing/\allowbreak{}SemigroupGlobal.lean, lines 133–150}
\begin{leancode}
\leanline{0}{theorem\ state\_integral\ (P\ :\ Parameters)\ (V\ :\ E\ P.dimension\ →\ ℝ)}
\leanline{4}{(hV\ :\ PotentialAssumptions\ P\ V)\ (t\ :\ ℝ)\ (ht\ :\ 0\ ≤\ t)\ :}
\leanline{4}{letI\ :=\ actualProbability\_isProbability\ P\ V\ hV}
\leanline{4}{state\ P\ V\ hV\ t\ -\ initial\ P\ V\ hV\ =}
\leanline{6}{-(∫\ s\ in\ 0..t,\ ULA.lpDrift\ (actualProbability\ P\ V)\ (gradient\ V)}
\leanline{8}{hV.gradientLipschitz\ (state\ P\ V\ hV\ s))\ +}
\leanline{6}{ULA.brownianLangevinNoise\ (actualProbability\ P\ V)\ (brownian\ P)}
\leanline{8}{(brownian\_preBrownian\ P\ V\ hV)\ 0\ t}
\end{leancode}
\leansource{\texttt{state\_\allowbreak{}eq\_\allowbreak{}actualSegment} \quad WassersteinMixing/\allowbreak{}SemigroupGlobal.lean, lines 153–169}
\begin{leancode}
\leanline{0}{theorem\ state\_eq\_actualSegment\ (P\ :\ Parameters)\ (V\ :\ E\ P.dimension\ →\ ℝ)}
\leanline{4}{(hV\ :\ PotentialAssumptions\ P\ V)\ (n\ :\ ℕ)\ (s\ :\ ℝ)\ (hs\ :\ s\ ∈\ Icc\ 0\ P.step)\ :}
\leanline{4}{state\ P\ V\ hV\ ((n\ :\ ℝ)\ *\ P.step\ +\ s)\ =\ actualSegment\ P\ V\ hV\ n\ s}
\end{leancode}
\leansource{\texttt{state\_\allowbreak{}grid} \quad WassersteinMixing/\allowbreak{}SemigroupGlobal.lean, lines 172–176}
\begin{leancode}
\leanline{0}{theorem\ state\_grid\ (P\ :\ Parameters)\ (V\ :\ E\ P.dimension\ →\ ℝ)}
\leanline{4}{(hV\ :\ PotentialAssumptions\ P\ V)\ (n\ :\ ℕ)\ :}
\leanline{4}{state\ P\ V\ hV\ ((n\ :\ ℝ)\ *\ P.step)\ =\ actualReference\ P\ V\ hV\ n}
\end{leancode}
\leanlabel{How the Lean proof follows the argument}
\texttt{state\_\allowbreak{}eq\_\allowbreak{}solution\allowbreak{}On} proves horizon independence by integral-solution uniqueness. The continuity lemmas use local equality with a longer fixed solution. \texttt{state\_\allowbreak{}integral}, \texttt{state\_\allowbreak{}eq\_\allowbreak{}actual\allowbreak{}Segment} and \texttt{state\_\allowbreak{}grid} carry that exact identification to the global and comparison representations.
\end{leanbox}


Earlier results used here: \cite[Section~3.13.2]{LangevinLean2026ULASC}; \cite[Section~3.25.1]{LangevinLean2026ULASC}; \cite[Section~3.43.4]{LangevinLean2026ULASC}; \cite[Section~3.44.5]{LangevinLean2026ULASC}.

\section{\texorpdfstring{Gibbs stationarity and the global observable contraction}{Gibbs stationarity and the global observable contraction}}


\textbf{Claim.} Every global state has law \(\pi\). Thus \(U_tf=f(Y_t)\) is an isometry, \(P_t=U_0^*U_t\) is a contraction, \(P_0=I\), and every orbit is strongly continuous for \(t\ge0\).



\textbf{Proof.} The actual Euler innovations have the independent coordinate and initialization laws required by the proved Brownian Gibbs-stationarity theorem. Apply it to the global finite-horizon solution to obtain the marginal law. Actual evaluation and compression then give the operator and norm identities. Continuity of the global state together with stationarity implies strong continuity of every observable evaluation; the fixed adjoint preserves it.

\begin{leanbox}{Lean: \texttt{scalar\_\allowbreak{}independent}, \texttt{initial\_\allowbreak{}independent}, \texttt{state\_\allowbreak{}law} and 5 more}
\leansource{\texttt{scalar\_\allowbreak{}independent} \quad WassersteinMixing/\allowbreak{}SemigroupGlobal.lean, lines 178–189}
\begin{leancode}
\leanline{0}{private\ theorem\ scalar\_independent\ (P\ :\ Parameters)\ (V\ :\ E\ P.dimension\ →\ ℝ)}
\leanline{4}{(hV\ :\ PotentialAssumptions\ P\ V)\ (r\ :\ ℝ)\ (hr\ :\ 0\ <\ r)\ :}
\leanline{4}{iIndepFun\ (fun\ p\ :\ ℕ\ ×\ Fin\ P.dimension\ ↦\ fun\ ω\ ↦}
\leanline{6}{ULA.Hilbert.brownianEulerNoise\ (actualProbability\ P\ V)\ (brownian\ P)}
\leanline{8}{(brownian\_preBrownian\ P\ V\ hV)\ r\ p.1\ ω\ p.2)\ (actualProbability\ P\ V)}
\end{leancode}
\leansource{\texttt{initial\_\allowbreak{}independent} \quad WassersteinMixing/\allowbreak{}SemigroupGlobal.lean, lines 191–202}
\begin{leancode}
\leanline{0}{private\ theorem\ initial\_independent\ (P\ :\ Parameters)\ (V\ :\ E\ P.dimension\ →\ ℝ)}
\leanline{4}{(hV\ :\ PotentialAssumptions\ P\ V)\ (r\ :\ ℝ)\ (\_hr\ :\ 0\ <\ r)\ :}
\leanline{4}{IndepFun\ (initial\ P\ V\ hV)\ (fun\ ω\ k\ ↦}
\leanline{6}{ULA.Hilbert.brownianEulerNoise\ (actualProbability\ P\ V)\ (brownian\ P)}
\leanline{8}{(brownian\_preBrownian\ P\ V\ hV)\ r\ k\ ω)\ (actualProbability\ P\ V)}
\end{leancode}
\leansource{\texttt{state\_\allowbreak{}law} \quad WassersteinMixing/\allowbreak{}SemigroupGlobal.lean, lines 206–217}
\begin{leancode}
\leanline{0}{theorem\ state\_law\ (P\ :\ Parameters)\ (V\ :\ E\ P.dimension\ →\ ℝ)}
\leanline{4}{(hV\ :\ PotentialAssumptions\ P\ V)\ (t\ :\ ℝ)\ :}
\leanline{4}{(actualProbability\ P\ V).map\ (state\ P\ V\ hV\ t)\ =\ ULA.gibbsMeasure\ V}
\end{leancode}
\leansource{\texttt{evaluation} \quad WassersteinMixing/\allowbreak{}SemigroupGlobal.lean, lines 223–226}
\begin{leancode}
\leanline{0}{def\ evaluation\ (P\ :\ Parameters)\ (V\ :\ E\ P.dimension\ →\ ℝ)\ (hV\ :\ PotentialAssumptions\ P\ V)}
\leanline{4}{(t\ :\ ℝ)\ :\ Lp\ F\ 2\ (ULA.gibbsMeasure\ V)\ →ₗᵢ[{\leanfallback 𝕜}]\ Lp\ F\ 2\ (actualProbability\ P\ V)\ :=}
\leanline{2}{StationaryObservables.evaluation\ (ULA.gibbsMeasure\ V)\ (actualProbability\ P\ V)}
\leanline{4}{(state\ P\ V\ hV\ t)\ (state\_law\ P\ V\ hV\ t)}
\end{leancode}
\leansource{\texttt{op} \quad WassersteinMixing/\allowbreak{}SemigroupGlobal.lean, lines 229–231}
\begin{leancode}
\leanline{0}{def\ op\ (P\ :\ Parameters)\ (V\ :\ E\ P.dimension\ →\ ℝ)\ (hV\ :\ PotentialAssumptions\ P\ V)\ (t\ :\ ℝ)\ :}
\leanline{4}{Lp\ F\ 2\ (ULA.gibbsMeasure\ V)\ →L[{\leanfallback 𝕜}]\ Lp\ F\ 2\ (ULA.gibbsMeasure\ V)\ :=}
\leanline{2}{ObservableTransition.compression\ (evaluation\ P\ V\ hV\ 0)\ (evaluation\ P\ V\ hV\ t)}
\end{leancode}
\leansource{\texttt{op\_\allowbreak{}zero} \quad WassersteinMixing/\allowbreak{}SemigroupGlobal.lean, lines 233–235}
\begin{leancode}
\leanline{0}{@[simp]\ theorem\ op\_zero\ (P\ :\ Parameters)\ (V\ :\ E\ P.dimension\ →\ ℝ)}
\leanline{4}{(hV\ :\ PotentialAssumptions\ P\ V)\ :\ op\ ({\leanfallback 𝕜}\ :=\ {\leanfallback 𝕜})\ (F\ :=\ F)\ P\ V\ hV\ 0\ =\ 1}
\end{leancode}
\leansource{\texttt{op\_\allowbreak{}norm\_\allowbreak{}le} \quad WassersteinMixing/\allowbreak{}SemigroupGlobal.lean, lines 237–240}
\begin{leancode}
\leanline{0}{theorem\ op\_norm\_le\ (P\ :\ Parameters)\ (V\ :\ E\ P.dimension\ →\ ℝ)}
\leanline{4}{(hV\ :\ PotentialAssumptions\ P\ V)\ (t\ :\ ℝ)\ (f\ :\ Lp\ F\ 2\ (ULA.gibbsMeasure\ V))\ :}
\leanline{4}{‖op\ ({\leanfallback 𝕜}\ :=\ {\leanfallback 𝕜})\ P\ V\ hV\ t\ f‖\ ≤\ ‖f‖}
\end{leancode}
\leansource{\texttt{continuousOn\_\allowbreak{}op} \quad WassersteinMixing/\allowbreak{}SemigroupGlobal.lean, lines 244–254}
\begin{leancode}
\leanline{0}{theorem\ continuousOn\_op\ [NormedSpace\ ℝ\ F]\ (P\ :\ Parameters)\ (V\ :\ E\ P.dimension\ →\ ℝ)}
\leanline{4}{(hV\ :\ PotentialAssumptions\ P\ V)\ (f\ :\ Lp\ F\ 2\ (ULA.gibbsMeasure\ V))\ :}
\leanline{4}{ContinuousOn\ (fun\ t\ ↦\ op\ ({\leanfallback 𝕜}\ :=\ {\leanfallback 𝕜})\ P\ V\ hV\ t\ f)\ (Ici\ 0)}
\end{leancode}
\leanlabel{How the Lean proof follows the argument}
The two independence helpers specialize the actual Brownian product source. \texttt{state\_\allowbreak{}law} invokes the proved Langevin Gibbs-stationarity result. \texttt{evaluation} and \texttt{op} are actual pullback and compression constructions; \texttt{continuous\allowbreak{}On\_\allowbreak{}op} uses the stationary-observable continuity theorem.
\end{leanbox}


Earlier results used here: \cite[Section~3.26.3]{LangevinLean2026ULASC}; \cite[Section~3.29.7]{LangevinLean2026ULASC}; \cite[Section~3.32.13]{LangevinLean2026ULASC}; \cite[Section~3.32.15]{LangevinLean2026ULASC}; \cite[Section~3.32.17]{LangevinLean2026ULASC}; \cite[Section~3.42.7]{LangevinLean2026ULASC}; \cite[Section~3.8.17]{LangevinLean2026ULASC}; \cite[Section~3.8.19]{LangevinLean2026ULASC}; \cite[Section~3.8.20]{LangevinLean2026ULASC}; \cite[Section~3.8.22]{LangevinLean2026ULASC}.

\begin{thebibliography}{ULA-LSI}
\bibitem[ULA-SC]{LangevinLean2026ULASC}
LangevinLean Project. \emph{Unadjusted Langevin under Strong Convexity}. LL-01, 2026. Draft report, content revision d34ead06987c. \url{https://odakyildiz.github.io/LangevinLean/docs/pdf/UnadjustedLangevin/FullProof.pdf}.
\end{thebibliography}

\end{document}
